\documentclass[conference]{IEEEtran}
\usepackage{amsmath,amssymb}
\usepackage{algorithm}
\usepackage{algpseudocode}
\usepackage{hyperref}

\begin{document}

\title{Reconstructable Gate Grouping and Priority-Cut Node Sharing for
Reversible Pebbling of Boolean Networks}

\author{\IEEEauthorblockN{Coding Agent Summary Report}
\IEEEauthorblockA{Generated from \texttt{pebbling\_solver.py}}}

\maketitle

\begin{abstract}
We describe the current gate-grouping and node-sharing algorithms of a
Z3-based reversible-pebbling framework over Boolean logic DAGs. Building
on a prior topological, contiguous-range gate-grouping scheme, we
identify and fix a second structural defect -- a gate whose boundary
outputs mix a genuine primary output with a merely-shared intermediate
value can never release that intermediate's qubit, since primary-output
gates never uncompute. We introduce a group-splitting pass that
restores correctness while preserving the acyclicity guarantee, and we
correspondingly fix the qubit-reclamation policy so that gate-internal
(dirty) values are freed eagerly rather than deferred to an uncompute
event that may never occur. We further describe a priority-cut-style
node-sharing scheme for Reed-Muller (ANF) network synthesis, which
generalizes simple exact-prefix caching to genuine subset-cut reuse
across monomials and across output functions.
\end{abstract}

\section{Recap: Topological Contiguous-Range Grouping}

Let $G=(V,E)$ be a combinational Boolean network DAG with primary
inputs $V_{PI}$, primary outputs $O\subseteq V$, and node list
$v_1,\dots,v_n$ maintained in topological order. A gate grouping
partitions the non-input nodes into groups $G_1,\dots,G_m$, each
required to be a contiguous range $\{v_{\ell_i},\dots,v_{r_i}\}$ of
this order. As established previously, any contiguous range of a
topological order is automatically convex, which in turn guarantees the
induced gate-dependency graph $\widehat G$ is acyclic, since
$\widehat G$ then does nothing more than reflect $G$'s own topological
order.

For each group $G_i$, boundary computation classifies every owned node
$v\in G_i$ as an \emph{output} (if $v\in O$ or some fanout of $v$ lies
outside $G_i$) or \emph{internal} (otherwise), and collects
$\mathrm{inputs}(G_i)$ as the set of external fanins consumed by $G_i$.
Gate-level dependencies are derived directly from
$\mathrm{inputs}(G_i)$ against node ownership.

\section{The Mixed-PO-Group Defect}

\texttt{GatePebbleSolver} requires any gate owning at least one genuine
primary output node (a member of $O$) to remain pebbled -- i.e.\
$p_{g,K}=\mathrm{true}$ -- for the remainder of the schedule. Such a
gate therefore never executes an \texttt{uncompute\_gate} transition.

This becomes a correctness defect when a single group's output set
$\mathrm{outputs}(G_i)$ contains \emph{both} a genuine PO node and a
non-PO output node -- a value that is merely consumed by a node in a
different, later group (a common occurrence once cut sharing lets a
monomial subexpression feed more than one output cone). The non-PO
output's qubit is allocated when $G_i$ computes, but the ONLY place a
non-PO output's qubit is released is $G_i$'s own
\texttt{uncompute\_gate} event -- which, if $G_i$ also owns a PO, never
occurs. The qubit therefore leaks permanently, which manifests as a
failed end-of-schedule liveness assertion:
\begin{equation}
\exists\, q:\ q \in \mathrm{live}(K) \ \wedge\ \mathrm{owner}(q) \notin O.
\end{equation}

\subsection{Fix: Mixed-PO Group Splitting}

We introduce a post-boundary-computation pass,
\texttt{\_split\_mixed\_po\_groups}, enforcing the invariant
\begin{equation}
\forall i:\quad
\big(\mathrm{outputs}(G_i)\cap O \ne \emptyset\big)
\;\Rightarrow\;
\mathrm{outputs}(G_i)\subseteq O.
\end{equation}
That is, a group's outputs are either entirely genuine primary outputs,
or contain no primary output at all -- never a mixture.

For any group violating this invariant, let $j^\ast$ be the largest
index (in the group's own node order) of a non-PO output node. The
group's owned-node list is split at $j^\ast$: nodes
$v_{\ell_i},\dots,v_{j^\ast}$ become an earlier group, and the
remainder $v_{j^\ast+1},\dots,v_{r_i}$ become a later group. Because
every group is a contiguous range and a sub-range of a contiguous range
is itself contiguous, this operation preserves convexity -- and
therefore acyclicity of $\widehat G$ -- without any additional search.
Boundaries and the invariant check are then recomputed; this repeats to
a fixed point, guaranteed to terminate since each split strictly
increases the group count while the total node count is fixed.

\begin{algorithm}
\caption{Mixed-PO group splitting (fixed point)}
\label{alg:splitmixed}
\begin{algorithmic}[1]
\Require groups $G_1,\dots,G_m$ (contiguous ranges); PO set $O$
\Ensure no group's outputs mix a PO with a non-PO output
\Repeat
  \State recompute $\mathrm{outputs}(G_i)$, $\mathrm{internal}(G_i)$, $\mathrm{inputs}(G_i)$ for all $i$
  \State $changed \gets \mathrm{false}$
  \For{each group $G_i$}
    \State $P \gets \mathrm{outputs}(G_i)\cap O$; \quad $N \gets \mathrm{outputs}(G_i)\setminus O$
    \If{$P \ne \emptyset$ \textbf{and} $N \ne \emptyset$}
      \State $j^\ast \gets$ index of last node in $N$ (by group order)
      \State split $G_i$'s node list at $j^\ast$ into two contiguous groups
      \State $changed \gets \mathrm{true}$
    \EndIf
  \EndFor
\Until{$changed = \mathrm{false}$}
\end{algorithmic}
\end{algorithm}

A defensive check, \texttt{\_validate\_no\_mixed\_po\_groups}, re-scans
the final grouping and raises if the invariant is somehow still
violated, guarding against regressions in the splitting logic itself.

\section{Qubit Reclamation Policy}

A second, related defect existed independently in qubit assignment.
Internal (dirty) nodes of a group are, by definition, values consumed
only within their own owning gate's compute step -- never by anything
outside the group. An earlier revision freed such nodes only at their
owning gate's \texttt{uncompute\_gate} event, which -- for exactly the
same reason as above -- never fires for a PO-owning gate, silently
leaking every internal ancilla qubit that gate ever used.

The corrected policy frees internal nodes \emph{eagerly}, immediately
after the owning gate's own \texttt{compute\_gate} step, rather than
deferring to its (possibly nonexistent) uncompute event:
\begin{equation}
\forall v \in \mathrm{internal}(G_i):\quad
t_{\mathrm{free}}(v) = t_{\mathrm{compute}}(G_i),
\end{equation}
in contrast to output nodes, whose qubits are released at
$t_{\mathrm{uncompute}}(G_i)$ if and only if $v\notin O$, and never
released if $v\in O$. With the mixed-PO-group invariant from
Section~III in place, a group with any non-PO output is now guaranteed
to eventually execute \texttt{uncompute\_gate}, making this release
path reachable in all cases where it is required.

\section{Updated Constraint Summary}

The Z3 encodings themselves are unchanged from the prior revision;
only the classical (non-Z3) construction and post-processing of gate
groups is affected. For completeness, the gate-level clause families
remain:

\paragraph{Initial and final clauses.}
\begin{equation}
\bigwedge_{g} \bar p_{g,0}
\;\wedge\;
\bigwedge_{g:\ \mathrm{outputs}(g)\cap O\ne\emptyset} p_{g,K}
\;\wedge\;
\bigwedge_{g:\ \mathrm{outputs}(g)\cap O=\emptyset} \bar p_{g,K}.
\end{equation}

\paragraph{Move clauses.}
\begin{equation}
\bigwedge_{i=1}^{K}\;\bigwedge_{g}\;\bigwedge_{h\in\mathrm{depends\_on}(g)}
\Big(\big(p_{g,i}\oplus p_{g,i+1}\big) \rightarrow
\big(p_{h,i}\wedge p_{h,i+1}\big)\Big).
\end{equation}

\paragraph{Cardinality clauses.}
\begin{equation}
\bigwedge_{i}\Big(\sum_{g:\ \mathrm{counts}(g)} p_{g,i} \le P\Big),
\end{equation}
where $\mathrm{counts}(g)$ excludes groups whose owned nodes are all
XOR nodes. With the Section III fix, the final-clause partition above
is now well-defined without ambiguity: every group falls cleanly into
exactly one of the two final-condition classes, since no group can
straddle both by owning a mixture of PO and non-PO outputs.

\section{Priority-Cut Node Sharing for Reed-Muller Synthesis}

The prior Reed-Muller network builder shared AND-chain nodes only when
two monomials happened to share an exact variable-index \emph{prefix}
(e.g.\ $(0,1,2)$ and $(0,1,2,3)$), and only within a single fixed
extension order. We generalize this to a priority-cut-style scheme,
mirroring the cut-based technology mapping used in tools such as ABC.

Each monomial is treated as a target variable set
$M\subseteq\{0,\dots,n-1\}$ to be realized as a chain of 2-input AND
gates. Rather than only checking $M$ itself against a cache, we search
the cache of \emph{all} previously constructed subset-AND nodes for a
usable \emph{cut}:
\begin{equation}
C^\ast = \arg\max_{C \in \mathcal{S},\ C \subsetneq M}
\big(|C|,\ \mathrm{reuse}(C)\big),
\end{equation}
where $\mathcal S$ is the set of all cached subset keys, ties are
broken lexicographically by the pair, and $\mathrm{reuse}(C)$ counts
how many times $C$ has previously been selected as a cut (favoring
concentration of reuse onto already-popular subexpressions, a proxy for
priority-cut mapping's area/depth-driven cut selection). If no proper
subset is cached, a single variable is used as the seed cut.

The remaining variables $M\setminus C^\ast$ are then appended one at a
time as a chain of 2-input AND gates, with every newly created
intermediate subset registered in the cache under its own frozenset
key -- making it available as a cut candidate for any subsequent
monomial, whether from the same output function or a different one.
Monomials are processed largest-first (by $|M|$) so that large,
information-rich cuts are registered in the cache before smaller
monomials search it, which empirically improves sharing yield.

\begin{algorithm}
\caption{Priority-cut AND-node construction for a monomial}
\label{alg:pcutand}
\begin{algorithmic}[1]
\Require monomial variable set $M$; subset cache $\mathcal S$ (fs $\to$ Node)
\Ensure Node realizing $\bigwedge_{i\in M} x_i$, with cache updated
\If{$M = \emptyset$} \Return \texttt{const1} \EndIf
\If{$M \in \mathcal S$} \Return $\mathcal S[M]$ \EndIf
\State $C^\ast \gets \arg\max_{C\in\mathcal S,\ C\subsetneq M} (|C|, \mathrm{reuse}(C))$
\If{no such $C^\ast$ exists}
  \State $C^\ast \gets \{\min(M)\}$ \Comment{seed with a single variable}
\EndIf
\State $\mathrm{reuse}(C^\ast) \mathrel{+}= 1$
\State $cur \gets C^\ast$; $node \gets \mathcal S[C^\ast]$
\For{$v \in \mathrm{sorted}(M \setminus C^\ast)$}
  \State $next \gets cur \cup \{v\}$
  \If{$next \in \mathcal S$}
    \State $node \gets \mathcal S[next]$; $\mathrm{reuse}(next) \mathrel{+}= 1$
  \Else
    \State $node \gets \mathrm{AND}(node, x_v)$; \quad $\mathcal S[next] \gets node$
  \EndIf
  \State $cur \gets next$
\EndFor
\State \Return $node$
\end{algorithmic}
\end{algorithm}

Because cut matching is by exact subset (not merely a shared prefix),
this scheme captures sharing opportunities the earlier prefix-chain
cache missed entirely -- for instance, monomials $(0,1,3)$ and
$(0,2,3)$ share no common prefix but both may usefully reuse a cached
cut such as $\{0,3\}$ if it happens to have been constructed first.
XOR reduction of an output's monomial terms continues to use
unordered-pair identity caching, unchanged from the prior revision.

\section{Conclusion}

Two structural gaps in the gate-grouping and qubit-assignment pipeline
-- mixed-PO output groups and deferred-instead-of-eager dirty-qubit
reclamation -- have been closed, both without weakening the acyclicity
guarantee established by contiguous-range grouping, since the fix in
each case is a further, still-contiguous partition refinement rather
than a change to the grouping topology itself. Separately, Reed-Muller
network synthesis now supports genuine priority-cut-style subset
sharing across monomials and across output functions, extending gate
count reduction beyond what exact-prefix caching alone could capture.

\end{document}
